\section*{Conditional Expectation}

\subsection*{13.1 Expectation Conditioned on a Finite Partition}

As a motivating example, suppose we want to obtain the average height of people

in a country by taking uniform samples from the population. We define the set of

people in this country as \Omega , and the uniform distribution on \Omega as P We denote the

height of a person chosen from the population by X()The expectation E[X]

represents the average height of the population. However, we can also compute

the average height in the country by first determining the average height in each

province. Let E[X\vertprovince A ] denote the average height in a given province A .

The national average height can then be computed as a weighted average of the

provincial averages, with the weighting factors proportional to the population in the

provinces.

We can also partition the people in this country by the cities they live in. The

average height in a city is the conditional expectation E[X\vertcity ]The partition of

the country into cities is a more refined partition compared to the partition into

provinces. This example illustrates that the conditional expectation is closely related

to how we partition the sample space.

In this chapter, we first review the conditional expectation of a random variable

given an event. Then we define the conditional expectation conditioned on a

\sigma-algebra generated by a finite number of events. In the abstract definition of

conditional expectation, we will condition on a general \sigma-algebra. Finally, we will

discuss filtrations and martingales as applications of conditional expectation.

Throughout this chapter, we will work with a fixed probability space .(\Omega,\mathcal{F},P) .

In probability theory, we define the conditional probability of an event E given

another event A by P(A \cap E)/P(A) , provided that P(A) > 0 This quantity

represents the likelihood of the event E given that we already know that the event

A has occurred.

If the event A is fixed and the event E varies, we obtain a measure function whose

value is conditional probability.

\begin{defbox}{13.1}
\end{defbox}

Let A be an event in \mathcal{F}withP(A) /=0. We define a probability measure \muA by

\muA(E) \coloneqq P(A \cap E)

P(A) .

The expectation conditioned on the event A is defined by the integral with respect

to the measure \muA,

E[X\vertA] \coloneqq

\int

\Omega

Xd\mu A.

It is easy to verify that \muA is a probability measure on \mathcal{F}with\muA(A) = 1In the

calculations of E[X] and E[X\vertA], we have the same function X mapping from \Omega

to the real numbers. The difference is in the change of measure from P to. \muA.

The next theorem gives the relationship between the two expectation operators.

\begin{thmbox}{13.1}
\end{thmbox}

For a random variable X inL1(P) ,

E[X\vertA]= 1

P(A)

\int

A

XdP.

\textit{Proof} Suppose X is an indicator function . 1B for someB \in\mathcal{F}Then we have

E[1B\vertA] \coloneqq

\int

\Omega

1B d\muA = \muA(B) \coloneqq P(A \cap B)

P(A) .

On the other hand,

P(A)

\int

A

1B dP = 1

P(A)

\int

\Omega

1A 1B dP = 1

P(A)

\int

\Omega

1A\capB dP = P(A \cap B)

P(A) .

Therefore, the equality in the theorem holds for \mathcal{F}-measurable indicator functions.

By linearity, we can extend it to all \mathcal{F}-measurable simple functions.

The remaining tasks are to extend this to nonnegative functions and real-valued

functions. By the monotone convergence theorem, we can extend it to all non-

negative \mathcal{F}-measurable functions. Finally, by decomposing real-valued functions

into positive and negative parts, we can extend it to all \mathcal{F}-measurable real-valued

functions. The rest of the proof is mechanical and is omitted. \hfill $\square$

We next extend the definition of conditional expectation to conditional expecta-

tion given a finite partition, as we would like to have the flexibility to condition on

different sets A.

\begin{defbox}{13.2}
\end{defbox}

Consider a finite partition. A1,A2,...,A k of the sample space. \Omega Let. \mathcal{G}be the\sigma-

algebra generated by A1,A2,...,A k For any \mathcal{G}-measurable function in L1(P) ,

define the conditional expectation given \mathcal{G}by

E[X\vert\mathcal{G}] \coloneqq

k\sum

i=1

E[X\vertAi]1Ai .

A few important notes are in order:

The conditional expectation given \mathcal{G}is a random variable. We can write it as

E[X\vert\mathcal{G}]() to emphasize that the input of the function is . .

E[X\vert\mathcal{G}] is a simple function that is \mathcal{G}-measurable and constant on each Ai The

value ofE[X\vert\mathcal{G}]() for. \in Ai isE[X\vertAi].

If Ai =\emptyset or P(Ai) = 0, the conditional expectation E[X\vertAi] given Ai is not

defined, butE[X\vert\mathcal{G}] is still well-defined.

Since E[X\vert\mathcal{G}] is a random variable, a natural question is how to compute the

expectation.

\begin{thmbox}{13.2}
\end{thmbox}

For any eventC \in\mathcal{G},

\int

C

E[X\vert\mathcal{G}] dP =

\int

C

XdP. (13.1)

In particular , we have

E[E[X\vert\mathcal{G}]] = E[X]. (13.2)

\textit{Proof} Recall that the events A1 to Ak form a partition of \Omega and are mutually

disjoint. They are the atoms of the \sigma-algebra \mathcal{G}Any event C in \mathcal{G}is a union of

some sets among A1 to Ak Thus, we can find an index set I \subseteq{ 1, 2,...,k } such

thatC =\cup i\inI Ai .

By the definition of expectation for simple function, we have

\int

C

E[X\vert\mathcal{G}] dP =

k\sum

i=1

E[X\vertAi]\cdot P(C \cap Ai) =

\sum

i\inI

E[X\vertAi]P(Ai)

=

\sum

i\inI

( 1

P(Ai)

\int

Ai

XdP

)

P(Ai),

which can be simplified to

\sum

i\inI

\int

Ai

XdP =

\int

C

XdP.

This holds for any C \in\mathcal{G}By setting C = \Omega , we obtain (13.2)\hfill $\square$

The equality in (13.2) is saying that the overall average valueE[X] is a weighted

average of the local average values E[X\vertAi]'s. In other words, the value of E[X\vert\mathcal{G}]

can be thought of as a "smoothing" of the random variable X that takes the partition

A1,A 2,...,A k into account.
